%%%%%%%%%%%%%%%%%%%%%%%%%%%%%%%%%%%%%%%%%%%%%%%%%%%%%%%%%%%%
% 12.3 MODEL THEORY
% The Composition of Advanced Mathematics
%%%%%%%%%%%%%%%%%%%%%%%%%%%%%%%%%%%%%%%%%%%%%%%%%%%%%%%%%%%%

\section{Model Theory}
\label{sec:model-theory}

\prerequisite{
Formal logic, first-order languages, syntax and semantics, theories,
structures, quantifiers, axiomatic set theory, functions, relations,
cardinality, and methods of proof.
}

\learningobjectives{
By the end of this section, the reader should be able to:
\begin{itemize}
    \item explain the purpose of model theory;
    \item distinguish languages, theories, and structures;
    \item define models of first-order theories;
    \item understand elementary equivalence;
    \item recognize isomorphisms of structures;
    \item distinguish isomorphism from elementary equivalence;
    \item understand embeddings and substructures;
    \item define elementary embeddings;
    \item understand elementary substructures;
    \item apply the Tarski--Vaught criterion conceptually;
    \item understand definable sets;
    \item distinguish parameter-free and parameter-definable sets;
    \item recognize definable functions and relations;
    \item understand complete theories;
    \item distinguish theory completeness from logical completeness;
    \item understand types conceptually;
    \item recognize complete types;
    \item understand realization and omission of types;
    \item understand compactness;
    \item apply compactness to construct nonstandard models;
    \item understand finite satisfiability;
    \item recognize elementary diagrams conceptually;
    \item understand the Löwenheim--Skolem theorems;
    \item recognize downward and upward Löwenheim--Skolem;
    \item understand countable elementary substructures;
    \item recognize the Skolem paradox;
    \item distinguish internal and external cardinality;
    \item understand categoricity;
    \item distinguish finite and infinite categoricity;
    \item recognize complete categorical theories conceptually;
    \item understand quantifier elimination conceptually;
    \item recognize model-complete theories;
    \item understand prime and saturated models conceptually;
    \item recognize homogeneous structures;
    \item understand back-and-forth arguments;
    \item recognize Fraïssé-type ideas conceptually;
    \item understand ultraproducts conceptually;
    \item recognize Łoś's theorem;
    \item understand ultrafilters in model-theoretic constructions;
    \item recognize nonstandard analysis conceptually;
    \item understand theories of algebraic structures;
    \item recognize model theory of fields and ordered structures;
    \item understand stability conceptually;
    \item recognize classification-theoretic questions;
    \item understand definable closure and algebraic closure conceptually;
    \item recognize model-theoretic independence;
    \item understand o-minimality conceptually;
    \item recognize applications to algebra, geometry, number theory,
          analysis, combinatorics, and computer science;
    \item connect model theory with formal logic, set theory, proof theory,
          computability, and foundations of mathematics.
\end{itemize}
}

%%%%%%%%%%%%%%%%%%%%%%%%%%%%%%%%%%%%%%%%%%%%%%%%%%%%%%%%%%%%
\subsection{What Is Model Theory?}
\label{subsec:what-is-model-theory}
%%%%%%%%%%%%%%%%%%%%%%%%%%%%%%%%%%%%%%%%%%%%%%%%%%%%%%%%%%%%

Model theory studies the relationship between formal languages and the
mathematical structures that interpret them.

A typical situation consists of:

\[
\boxed{
\text{language}
}
\]

\[
\downarrow
\]

\[
\boxed{
\text{theory}
}
\]

\[
\downarrow
\]

\[
\boxed{
\text{class of models}.
}
\]

The language specifies what can be expressed.

The theory specifies which sentences are assumed.

The models are the structures in which those sentences are true.

%%%%%%%%%%%%%%%%%%%%%%%%%%%%%%%%%%%%%%%%%%%%%%%%%%%%%%%%%%%%
\subsection{Languages, Theories, and Models}
\label{subsec:languages-theories-models}
%%%%%%%%%%%%%%%%%%%%%%%%%%%%%%%%%%%%%%%%%%%%%%%%%%%%%%%%%%%%

A first-order language \(L\) contains logical symbols together with
nonlogical symbols such as:

\[
\boxed{
\text{constant symbols},
}
\]

\[
\boxed{
\text{function symbols},
}
\]

and

\[
\boxed{
\text{relation symbols}.
}
\]

A theory \(T\) is a set of \(L\)-sentences.

An \(L\)-structure \(\mathcal{M}\) is a model of \(T\) when

\[
\boxed{
\mathcal{M}\models T.
}
\]

This means

\[
\mathcal{M}\models\varphi
\]

for every

\[
\varphi\in T.
\]

%%%%%%%%%%%%%%%%%%%%%%%%%%%%%%%%%%%%%%%%%%%%%%%%%%%%%%%%%%%%
\subsection{Classes of Models}
\label{subsec:classes-of-models}
%%%%%%%%%%%%%%%%%%%%%%%%%%%%%%%%%%%%%%%%%%%%%%%%%%%%%%%%%%%%

For a theory \(T\), write

\[
\boxed{
\operatorname{Mod}(T)
}
\]

for the class of all models of \(T\).

Thus

\[
\boxed{
\operatorname{Mod}(T)
=
\{
\mathcal{M}:
\mathcal{M}\models T
\}.
}
\]

Model theory studies the structure of this class and the relationships
among its members.

%%%%%%%%%%%%%%%%%%%%%%%%%%%%%%%%%%%%%%%%%%%%%%%%%%%%%%%%%%%%
\subsection{Examples of First-Order Theories}
\label{subsec:first-order-theory-examples}
%%%%%%%%%%%%%%%%%%%%%%%%%%%%%%%%%%%%%%%%%%%%%%%%%%%%%%%%%%%%

Important examples include:

\[
\boxed{
\text{group theory},
}
\]

\[
\boxed{
\text{field theory},
}
\]

\[
\boxed{
\text{linear-order theory},
}
\]

\[
\boxed{
\text{graph theory},
}
\]

and

\[
\boxed{
\text{arithmetic}.
}
\]

Each theory describes a class of mathematical structures using first-order
sentences.

%%%%%%%%%%%%%%%%%%%%%%%%%%%%%%%%%%%%%%%%%%%%%%%%%%%%%%%%%%%%
\subsection{Theory of a Structure}
\label{subsec:theory-of-a-structure}
%%%%%%%%%%%%%%%%%%%%%%%%%%%%%%%%%%%%%%%%%%%%%%%%%%%%%%%%%%%%

For an \(L\)-structure \(\mathcal{M}\), define

\[
\boxed{
\operatorname{Th}(\mathcal{M})
=
\{
\varphi:
\mathcal{M}\models\varphi
\},
}
\]

where \(\varphi\) ranges over \(L\)-sentences.

Thus

\[
\operatorname{Th}(\mathcal{M})
\]

contains every first-order sentence true in \(\mathcal{M}\).

%%%%%%%%%%%%%%%%%%%%%%%%%%%%%%%%%%%%%%%%%%%%%%%%%%%%%%%%%%%%
\subsection{Elementary Equivalence}
\label{subsec:elementary-equivalence}
%%%%%%%%%%%%%%%%%%%%%%%%%%%%%%%%%%%%%%%%%%%%%%%%%%%%%%%%%%%%

\begin{definitionbox}{Elementary Equivalence}
Two \(L\)-structures \(\mathcal{M}\) and \(\mathcal{N}\) are elementarily
equivalent if they satisfy exactly the same first-order sentences.

Write

\[
\boxed{
\mathcal{M}
\equiv
\mathcal{N}.
}
\]
\end{definitionbox}

Equivalently,

\[
\boxed{
\operatorname{Th}(\mathcal{M})
=
\operatorname{Th}(\mathcal{N}).
}
\]

%%%%%%%%%%%%%%%%%%%%%%%%%%%%%%%%%%%%%%%%%%%%%%%%%%%%%%%%%%%%
\subsection{Isomorphism}
\label{subsec:model-theoretic-isomorphism}
%%%%%%%%%%%%%%%%%%%%%%%%%%%%%%%%%%%%%%%%%%%%%%%%%%%%%%%%%%%%

An isomorphism between structures is a bijection preserving all interpreted
constants, functions, and relations.

If

\[
f:\mathcal{M}\to\mathcal{N}
\]

is an isomorphism, then

\[
\boxed{
\mathcal{M}\cong\mathcal{N}.
}
\]

Isomorphic structures have the same mathematical structure up to renaming
of elements.

%%%%%%%%%%%%%%%%%%%%%%%%%%%%%%%%%%%%%%%%%%%%%%%%%%%%%%%%%%%%
\subsection{Isomorphism Implies Elementary Equivalence}
\label{subsec:isomorphism-elementary-equivalence}
%%%%%%%%%%%%%%%%%%%%%%%%%%%%%%%%%%%%%%%%%%%%%%%%%%%%%%%%%%%%

If

\[
\mathcal{M}\cong\mathcal{N},
\]

then every first-order sentence has the same truth value in both
structures.

Therefore,

\[
\boxed{
\mathcal{M}\cong\mathcal{N}
\Longrightarrow
\mathcal{M}\equiv\mathcal{N}.
}
\]

The converse is generally false.

%%%%%%%%%%%%%%%%%%%%%%%%%%%%%%%%%%%%%%%%%%%%%%%%%%%%%%%%%%%%
\subsection{Elementary Equivalence Is Weaker Than Isomorphism}
\label{subsec:elementary-equivalence-weaker}
%%%%%%%%%%%%%%%%%%%%%%%%%%%%%%%%%%%%%%%%%%%%%%%%%%%%%%%%%%%%

Two structures can satisfy exactly the same first-order sentences without
being isomorphic.

In particular, structures of different infinite cardinalities may be
elementarily equivalent.

This phenomenon is central to the Löwenheim--Skolem theorems.

%%%%%%%%%%%%%%%%%%%%%%%%%%%%%%%%%%%%%%%%%%%%%%%%%%%%%%%%%%%%
\subsection{Worked Example: Isomorphic Structures}
\label{subsec:worked-isomorphic-structures}
%%%%%%%%%%%%%%%%%%%%%%%%%%%%%%%%%%%%%%%%%%%%%%%%%%%%%%%%%%%%

\begin{workedexamplebox}{Two Copies of the Integers}

Let

\[
\mathcal{M}
=
(\mathbb{Z},+)
\]

and let

\[
\mathcal{N}
=
(2\mathbb{Z},+).
\]

Define

\[
f:\mathbb{Z}\to2\mathbb{Z}
\]

by

\[
f(n)=2n.
\]

This map is bijective and satisfies

\[
f(m+n)
=
2(m+n)
=
2m+2n
=
f(m)+f(n).
\]

Therefore,

\begin{answerbox}
\[
\boxed{
(\mathbb{Z},+)
\cong
(2\mathbb{Z},+).
}
\]
\end{answerbox}

Hence the two structures are also elementarily equivalent.

\end{workedexamplebox}

%%%%%%%%%%%%%%%%%%%%%%%%%%%%%%%%%%%%%%%%%%%%%%%%%%%%%%%%%%%%
\subsection{Embeddings}
\label{subsec:model-theoretic-embeddings}
%%%%%%%%%%%%%%%%%%%%%%%%%%%%%%%%%%%%%%%%%%%%%%%%%%%%%%%%%%%%

An embedding is an injective map preserving the structure.

For relational and functional languages, an embedding preserves the
interpretations of the symbols in the appropriate sense.

Conceptually,

\[
\boxed{
\mathcal{M}
\hookrightarrow
\mathcal{N}.
}
\]

An embedding identifies \(\mathcal{M}\) with a substructure of
\(\mathcal{N}\).

%%%%%%%%%%%%%%%%%%%%%%%%%%%%%%%%%%%%%%%%%%%%%%%%%%%%%%%%%%%%
\subsection{Substructures}
\label{subsec:model-theoretic-substructures}
%%%%%%%%%%%%%%%%%%%%%%%%%%%%%%%%%%%%%%%%%%%%%%%%%%%%%%%%%%%%

Suppose \(\mathcal{M}\) and \(\mathcal{N}\) are \(L\)-structures with

\[
M\subseteq N.
\]

Then \(\mathcal{M}\) is a substructure of \(\mathcal{N}\), written

\[
\boxed{
\mathcal{M}
\subseteq
\mathcal{N},
}
\]

when:

\begin{itemize}
    \item constant interpretations agree;
    \item \(M\) is closed under interpreted functions;
    \item relations on \(M\) are inherited from \(\mathcal{N}\).
\end{itemize}

%%%%%%%%%%%%%%%%%%%%%%%%%%%%%%%%%%%%%%%%%%%%%%%%%%%%%%%%%%%%
\subsection{Substructure Does Not Mean Elementary Substructure}
\label{subsec:substructure-not-elementary}
%%%%%%%%%%%%%%%%%%%%%%%%%%%%%%%%%%%%%%%%%%%%%%%%%%%%%%%%%%%%

A substructure may fail to preserve truth of formulas involving
quantifiers.

For example, a larger structure may contain witnesses to an existential
formula that do not lie in the smaller structure.

Thus

\[
\boxed{
\mathcal{M}\subseteq\mathcal{N}
}
\]

does not generally imply

\[
\boxed{
\mathcal{M}\preccurlyeq\mathcal{N}.
}
\]

%%%%%%%%%%%%%%%%%%%%%%%%%%%%%%%%%%%%%%%%%%%%%%%%%%%%%%%%%%%%
\subsection{Elementary Substructures}
\label{subsec:elementary-substructures}
%%%%%%%%%%%%%%%%%%%%%%%%%%%%%%%%%%%%%%%%%%%%%%%%%%%%%%%%%%%%

\begin{definitionbox}{Elementary Substructure}
Let

\[
\mathcal{M}
\subseteq
\mathcal{N}.
\]

We say \(\mathcal{M}\) is an elementary substructure of \(\mathcal{N}\),
written

\[
\boxed{
\mathcal{M}
\preccurlyeq
\mathcal{N},
}
\]

if for every first-order formula

\[
\varphi(x_1,\ldots,x_n)
\]

and every

\[
a_1,\ldots,a_n\in M,
\]

we have

\[
\boxed{
\mathcal{M}
\models
\varphi(a_1,\ldots,a_n)
\iff
\mathcal{N}
\models
\varphi(a_1,\ldots,a_n).
}
\]
\end{definitionbox}

%%%%%%%%%%%%%%%%%%%%%%%%%%%%%%%%%%%%%%%%%%%%%%%%%%%%%%%%%%%%
\subsection{Elementary Embeddings}
\label{subsec:elementary-embeddings}
%%%%%%%%%%%%%%%%%%%%%%%%%%%%%%%%%%%%%%%%%%%%%%%%%%%%%%%%%%%%

An embedding

\[
f:\mathcal{M}\to\mathcal{N}
\]

is elementary if it preserves truth of every first-order formula with
parameters.

Thus

\[
\boxed{
\mathcal{M}
\models
\varphi(a_1,\ldots,a_n)
}
\]

if and only if

\[
\boxed{
\mathcal{N}
\models
\varphi
(
f(a_1),\ldots,f(a_n)
).
}
\]

Isomorphisms are automatically elementary embeddings.

%%%%%%%%%%%%%%%%%%%%%%%%%%%%%%%%%%%%%%%%%%%%%%%%%%%%%%%%%%%%
\subsection{Tarski--Vaught Criterion}
\label{subsec:tarski-vaught-criterion}
%%%%%%%%%%%%%%%%%%%%%%%%%%%%%%%%%%%%%%%%%%%%%%%%%%%%%%%%%%%%

The Tarski--Vaught criterion gives a practical test for elementary
substructures.

Suppose

\[
\mathcal{M}
\subseteq
\mathcal{N}.
\]

Then \(\mathcal{M}\preccurlyeq\mathcal{N}\) if whenever

\[
\mathcal{N}
\models
\exists x\,\varphi(x,\mathbf{a}),
\]

with parameters

\[
\mathbf{a}\in M,
\]

there exists

\[
b\in M
\]

such that

\[
\boxed{
\mathcal{N}
\models
\varphi(b,\mathbf{a}).
}
\]

Thus every definable existential witness over \(M\) can already be found
inside \(M\).

%%%%%%%%%%%%%%%%%%%%%%%%%%%%%%%%%%%%%%%%%%%%%%%%%%%%%%%%%%%%
\subsection{Definable Sets}
\label{subsec:model-theoretic-definable-sets}
%%%%%%%%%%%%%%%%%%%%%%%%%%%%%%%%%%%%%%%%%%%%%%%%%%%%%%%%%%%%

A subset

\[
X\subseteq M^n
\]

is definable in \(\mathcal{M}\) if there exists a first-order formula

\[
\varphi(x_1,\ldots,x_n)
\]

such that

\[
\boxed{
X
=
\{
\mathbf{a}\in M^n:
\mathcal{M}\models\varphi(\mathbf{a})
\}.
}
\]

Definability is one of the central ideas of model theory.

%%%%%%%%%%%%%%%%%%%%%%%%%%%%%%%%%%%%%%%%%%%%%%%%%%%%%%%%%%%%
\subsection{Definability With Parameters}
\label{subsec:definability-with-parameters}
%%%%%%%%%%%%%%%%%%%%%%%%%%%%%%%%%%%%%%%%%%%%%%%%%%%%%%%%%%%%

A set may be definable using fixed elements of the structure as parameters.

For example,

\[
X
=
\{
x\in M:
\mathcal{M}\models
\varphi(x,a)
\}
\]

is definable with parameter \(a\).

Without parameters, one requires a formula containing no named elements of
\(M\).

%%%%%%%%%%%%%%%%%%%%%%%%%%%%%%%%%%%%%%%%%%%%%%%%%%%%%%%%%%%%
\subsection{Worked Example: Definable Subset of the Reals}
\label{subsec:worked-definable-set}
%%%%%%%%%%%%%%%%%%%%%%%%%%%%%%%%%%%%%%%%%%%%%%%%%%%%%%%%%%%%

\begin{workedexamplebox}{Positive Elements}

Consider the ordered field

\[
\mathcal{R}
=
(\mathbb{R},+,\cdot,<,0,1).
\]

The positive real numbers are defined by

\[
\varphi(x)
:
0<x.
\]

Therefore,

\begin{answerbox}
\[
\boxed{
\mathbb{R}_{>0}
=
\{
x\in\mathbb{R}:
\mathcal{R}\models0<x
\}.
}
\]
\end{answerbox}

Thus positivity is first-order definable in the ordered-field language.

\end{workedexamplebox}

%%%%%%%%%%%%%%%%%%%%%%%%%%%%%%%%%%%%%%%%%%%%%%%%%%%%%%%%%%%%
\subsection{Definable Functions}
\label{subsec:model-theoretic-definable-functions}
%%%%%%%%%%%%%%%%%%%%%%%%%%%%%%%%%%%%%%%%%%%%%%%%%%%%%%%%%%%%

A function

\[
f:M^n\to M
\]

is definable if its graph

\[
\boxed{
\Gamma_f
=
\{
(\mathbf{x},y):
y=f(\mathbf{x})
\}
}
\]

is a definable subset of

\[
M^{n+1}.
\]

Thus model theory treats functions through definability of their graphs.

%%%%%%%%%%%%%%%%%%%%%%%%%%%%%%%%%%%%%%%%%%%%%%%%%%%%%%%%%%%%
\subsection{Definable Closure}
\label{subsec:definable-closure}
%%%%%%%%%%%%%%%%%%%%%%%%%%%%%%%%%%%%%%%%%%%%%%%%%%%%%%%%%%%%

Given

\[
A\subseteq M,
\]

the definable closure of \(A\), denoted

\[
\boxed{
\operatorname{dcl}(A),
}
\]

consists of elements uniquely definable using parameters from \(A\).

Conceptually,

\[
b\in\operatorname{dcl}(A)
\]

if some formula over \(A\) uniquely identifies \(b\).

%%%%%%%%%%%%%%%%%%%%%%%%%%%%%%%%%%%%%%%%%%%%%%%%%%%%%%%%%%%%
\subsection{Algebraic Closure in Model Theory}
\label{subsec:model-theoretic-algebraic-closure}
%%%%%%%%%%%%%%%%%%%%%%%%%%%%%%%%%%%%%%%%%%%%%%%%%%%%%%%%%%%%

The model-theoretic algebraic closure of \(A\), denoted

\[
\boxed{
\operatorname{acl}(A),
}
\]

consists of elements lying in finite sets definable over \(A\).

Thus

\[
\boxed{
\operatorname{dcl}(A)
\subseteq
\operatorname{acl}(A).
}
\]

In theories of fields, this notion is closely related to ordinary algebraic
closure.

%%%%%%%%%%%%%%%%%%%%%%%%%%%%%%%%%%%%%%%%%%%%%%%%%%%%%%%%%%%%
\subsection{Complete Theories}
\label{subsec:complete-theories-model-theory}
%%%%%%%%%%%%%%%%%%%%%%%%%%%%%%%%%%%%%%%%%%%%%%%%%%%%%%%%%%%%

\begin{definitionbox}{Complete Theory}
A theory \(T\) is complete if for every sentence \(\varphi\) in its
language,

\[
\boxed{
T\models\varphi
}
\]

or

\[
\boxed{
T\models\neg\varphi.
}
\]
\end{definitionbox}

Equivalently, all models of \(T\) satisfy exactly the same sentences.

%%%%%%%%%%%%%%%%%%%%%%%%%%%%%%%%%%%%%%%%%%%%%%%%%%%%%%%%%%%%
\subsection{Completeness of a Theory Versus Logical Completeness}
\label{subsec:theory-vs-logical-completeness}
%%%%%%%%%%%%%%%%%%%%%%%%%%%%%%%%%%%%%%%%%%%%%%%%%%%%%%%%%%%%

These two notions must be distinguished.

A proof system is logically complete when

\[
\boxed{
T\models\varphi
\Longrightarrow
T\vdash\varphi.
}
\]

A theory \(T\) is complete when it decides every sentence:

\[
\boxed{
T\models\varphi
\quad
\text{or}
\quad
T\models\neg\varphi.
}
\]

The first is a property of a proof calculus.

The second is a property of a theory.

%%%%%%%%%%%%%%%%%%%%%%%%%%%%%%%%%%%%%%%%%%%%%%%%%%%%%%%%%%%%
\subsection{Complete Theory of a Structure}
\label{subsec:complete-theory-structure}
%%%%%%%%%%%%%%%%%%%%%%%%%%%%%%%%%%%%%%%%%%%%%%%%%%%%%%%%%%%%

For every structure \(\mathcal{M}\),

\[
\operatorname{Th}(\mathcal{M})
\]

is complete.

For any sentence \(\varphi\),

either

\[
\mathcal{M}\models\varphi
\]

or

\[
\mathcal{M}\models\neg\varphi.
\]

Therefore,

\[
\boxed{
\operatorname{Th}(\mathcal{M})
}
\]

decides every sentence in the language.

%%%%%%%%%%%%%%%%%%%%%%%%%%%%%%%%%%%%%%%%%%%%%%%%%%%%%%%%%%%%
\subsection{Types}
\label{subsec:model-theoretic-types}
%%%%%%%%%%%%%%%%%%%%%%%%%%%%%%%%%%%%%%%%%%%%%%%%%%%%%%%%%%%%

A type describes a possible collection of properties of an element or
tuple.

A set of formulas

\[
p(x)
\]

over theory \(T\) is a type when the formulas are jointly consistent with
\(T\).

Conceptually,

\[
\boxed{
p(x)
=
\{
\varphi_1(x),
\varphi_2(x),
\ldots
\}.
}
\]

A type can be viewed as a complete or partial description of a potential
element.

%%%%%%%%%%%%%%%%%%%%%%%%%%%%%%%%%%%%%%%%%%%%%%%%%%%%%%%%%%%%
\subsection{Partial Types}
\label{subsec:partial-types}
%%%%%%%%%%%%%%%%%%%%%%%%%%%%%%%%%%%%%%%%%%%%%%%%%%%%%%%%%%%%

A partial type is any consistent collection of formulas in specified free
variables.

For example, in an ordered structure one might consider

\[
\boxed{
p(x)
=
\{
x>0,
x>1,
x>2,
x>3,
\ldots
\}.
}
\]

This describes an element larger than every standard natural number named
in the language.

%%%%%%%%%%%%%%%%%%%%%%%%%%%%%%%%%%%%%%%%%%%%%%%%%%%%%%%%%%%%
\subsection{Complete Types}
\label{subsec:complete-types}
%%%%%%%%%%%%%%%%%%%%%%%%%%%%%%%%%%%%%%%%%%%%%%%%%%%%%%%%%%%%

A complete type decides every formula in the relevant variables.

For every formula

\[
\varphi(x),
\]

a complete type \(p(x)\) contains exactly one of

\[
\boxed{
\varphi(x)
}
\]

or

\[
\boxed{
\neg\varphi(x).
}
\]

Thus a complete type gives a maximally detailed consistent first-order
description.

%%%%%%%%%%%%%%%%%%%%%%%%%%%%%%%%%%%%%%%%%%%%%%%%%%%%%%%%%%%%
\subsection{Realization of a Type}
\label{subsec:realization-of-types}
%%%%%%%%%%%%%%%%%%%%%%%%%%%%%%%%%%%%%%%%%%%%%%%%%%%%%%%%%%%%

An element

\[
a\in M
\]

realizes a type \(p(x)\) if

\[
\boxed{
\mathcal{M}\models\varphi(a)
}
\]

for every

\[
\varphi(x)\in p.
\]

If no element realizes \(p\), the type is omitted in \(\mathcal{M}\).

%%%%%%%%%%%%%%%%%%%%%%%%%%%%%%%%%%%%%%%%%%%%%%%%%%%%%%%%%%%%
\subsection{Worked Example: A Type Not Realized in \(\mathbb{N}\)}
\label{subsec:worked-unrealized-type}
%%%%%%%%%%%%%%%%%%%%%%%%%%%%%%%%%%%%%%%%%%%%%%%%%%%%%%%%%%%%

\begin{workedexamplebox}{Larger Than Every Standard Numeral}

In arithmetic, consider

\[
p(x)
=
\{
x>0,
x>1,
x>2,
x>3,\ldots
\}.
\]

Every finite subset is realized in \(\mathbb{N}\).

For example,

\[
\{
x>0,
x>1,
\ldots,
x>100
\}
\]

is realized by \(101\).

But no standard natural number is larger than every natural number.

Thus

\begin{answerbox}
\[
\boxed{
p(x)
\text{ is finitely satisfiable in }\mathbb{N}
\text{ but not realized in }\mathbb{N}.
}
\]
\end{answerbox}

Compactness can produce an elementary extension in which such a type is
realized.

\end{workedexamplebox}

%%%%%%%%%%%%%%%%%%%%%%%%%%%%%%%%%%%%%%%%%%%%%%%%%%%%%%%%%%%%
\subsection{Compactness Theorem}
\label{subsec:model-theory-compactness}
%%%%%%%%%%%%%%%%%%%%%%%%%%%%%%%%%%%%%%%%%%%%%%%%%%%%%%%%%%%%

\begin{theorem}
A first-order theory \(T\) has a model if and only if every finite subset

\[
T_0\subseteq T
\]

has a model.
\end{theorem}

Thus

\[
\boxed{
T
\text{ satisfiable}
\iff
\text{every finite subset of }T
\text{ satisfiable}.
}
\]

Compactness is one of the central tools of model theory.

%%%%%%%%%%%%%%%%%%%%%%%%%%%%%%%%%%%%%%%%%%%%%%%%%%%%%%%%%%%%
\subsection{Semantic Consequence Form of Compactness}
\label{subsec:compactness-consequence-form}
%%%%%%%%%%%%%%%%%%%%%%%%%%%%%%%%%%%%%%%%%%%%%%%%%%%%%%%%%%%%

Compactness can also be stated as follows.

If

\[
T\models\varphi,
\]

then there exists a finite subset

\[
T_0\subseteq T
\]

such that

\[
\boxed{
T_0\models\varphi.
}
\]

Thus every first-order semantic consequence depends on only finitely many
assumptions.

%%%%%%%%%%%%%%%%%%%%%%%%%%%%%%%%%%%%%%%%%%%%%%%%%%%%%%%%%%%%
\subsection{Compactness and Infinite Models}
\label{subsec:compactness-infinite-models}
%%%%%%%%%%%%%%%%%%%%%%%%%%%%%%%%%%%%%%%%%%%%%%%%%%%%%%%%%%%%

Suppose a theory \(T\) has arbitrarily large finite models.

Introduce sentences

\[
\sigma_n
\]

asserting that there are at least \(n\) elements.

Then every finite subset of

\[
T
\cup
\{
\sigma_1,\sigma_2,\ldots
\}
\]

has a model.

Compactness gives an infinite model.

%%%%%%%%%%%%%%%%%%%%%%%%%%%%%%%%%%%%%%%%%%%%%%%%%%%%%%%%%%%%
\subsection{Worked Example: Graphs With Infinite Models}
\label{subsec:worked-compactness-graph}
%%%%%%%%%%%%%%%%%%%%%%%%%%%%%%%%%%%%%%%%%%%%%%%%%%%%%%%%%%%%

\begin{workedexamplebox}{Producing an Infinite Graph}

Suppose a graph theory \(T\) has finite models of arbitrarily large size.

Add sentences

\[
\sigma_n
=
\text{``there are at least \(n\) vertices.''}
\]

for every natural number \(n\).

Any finite collection of these sentences requires only a sufficiently large
finite graph.

Therefore every finite subset is satisfiable.

By compactness,

\begin{answerbox}
\[
\boxed{
T
\text{ has an infinite model.}
}
\]
\end{answerbox}

\end{workedexamplebox}

%%%%%%%%%%%%%%%%%%%%%%%%%%%%%%%%%%%%%%%%%%%%%%%%%%%%%%%%%%%%
\subsection{Compactness and Nonstandard Arithmetic}
\label{subsec:compactness-nonstandard-arithmetic}
%%%%%%%%%%%%%%%%%%%%%%%%%%%%%%%%%%%%%%%%%%%%%%%%%%%%%%%%%%%%

Let \(T\) be a first-order theory of arithmetic.

Extend the language with a new constant symbol \(c\).

Add sentences

\[
\boxed{
c>0,
\quad
c>1,
\quad
c>2,
\quad
\ldots.
}
\]

Every finite subset is satisfiable in the standard natural numbers by
interpreting \(c\) as a sufficiently large number.

Compactness therefore yields a model containing an element larger than
every standard numeral.

Such an element is nonstandard.

%%%%%%%%%%%%%%%%%%%%%%%%%%%%%%%%%%%%%%%%%%%%%%%%%%%%%%%%%%%%
\subsection{Nonstandard Models}
\label{subsec:nonstandard-models}
%%%%%%%%%%%%%%%%%%%%%%%%%%%%%%%%%%%%%%%%%%%%%%%%%%%%%%%%%%%%

A nonstandard model of arithmetic satisfies the same first-order axioms as
the intended natural-number structure but contains elements not
corresponding to ordinary finite natural numbers.

Thus first-order axioms for arithmetic may have models

\[
\boxed{
\mathcal{M}
\not\cong
\mathbb{N}.
}
\]

This demonstrates the difference between axiomatic description and
categorical determination.

%%%%%%%%%%%%%%%%%%%%%%%%%%%%%%%%%%%%%%%%%%%%%%%%%%%%%%%%%%%%
\subsection{Elementary Extensions}
\label{subsec:elementary-extensions}
%%%%%%%%%%%%%%%%%%%%%%%%%%%%%%%%%%%%%%%%%%%%%%%%%%%%%%%%%%%%

If

\[
\mathcal{M}
\preccurlyeq
\mathcal{N},
\]

then \(\mathcal{N}\) is an elementary extension of \(\mathcal{M}\).

Every first-order statement about elements of \(M\) has the same truth
value in the larger structure.

However, \(\mathcal{N}\) may contain many new elements.

%%%%%%%%%%%%%%%%%%%%%%%%%%%%%%%%%%%%%%%%%%%%%%%%%%%%%%%%%%%%
\subsection{Elementary Diagrams}
\label{subsec:elementary-diagrams}
%%%%%%%%%%%%%%%%%%%%%%%%%%%%%%%%%%%%%%%%%%%%%%%%%%%%%%%%%%%%

To study elementary extensions, one may expand the language with a constant
symbol

\[
c_a
\]

for every

\[
a\in M.
\]

The elementary diagram of \(\mathcal{M}\) contains all formulas with these
constants that are true in \(\mathcal{M}\).

Models of the elementary diagram correspond naturally to elementary
extensions of \(\mathcal{M}\).

%%%%%%%%%%%%%%%%%%%%%%%%%%%%%%%%%%%%%%%%%%%%%%%%%%%%%%%%%%%%
\subsection{Downward Löwenheim--Skolem Theorem}
\label{subsec:downward-lowenheim-skolem}
%%%%%%%%%%%%%%%%%%%%%%%%%%%%%%%%%%%%%%%%%%%%%%%%%%%%%%%%%%%%

One standard form states that if a first-order structure \(\mathcal{M}\)
is infinite and the language has cardinality at most \(\kappa\), then for
suitable infinite \(\kappa\) and parameter sets \(A\subseteq M\) of size at
most \(\kappa\), there exists an elementary substructure

\[
\mathcal{N}
\preccurlyeq
\mathcal{M}
\]

with

\[
A\subseteq N
\]

and

\[
\boxed{
|N|
\leq
\kappa.
}
\]

For a countable language and countable parameter set, one obtains a
countable elementary substructure.

%%%%%%%%%%%%%%%%%%%%%%%%%%%%%%%%%%%%%%%%%%%%%%%%%%%%%%%%%%%%
\subsection{Upward Löwenheim--Skolem Theorem}
\label{subsec:upward-lowenheim-skolem}
%%%%%%%%%%%%%%%%%%%%%%%%%%%%%%%%%%%%%%%%%%%%%%%%%%%%%%%%%%%%

If a first-order theory has an infinite model, then under standard
cardinality assumptions it has models in arbitrarily large infinite
cardinalities.

Thus an infinite first-order theory generally does not determine one
infinite model size.

This has major consequences for categoricity.

%%%%%%%%%%%%%%%%%%%%%%%%%%%%%%%%%%%%%%%%%%%%%%%%%%%%%%%%%%%%
\subsection{The Löwenheim--Skolem Phenomenon}
\label{subsec:lowenheim-skolem-phenomenon}
%%%%%%%%%%%%%%%%%%%%%%%%%%%%%%%%%%%%%%%%%%%%%%%%%%%%%%%%%%%%

Together, the downward and upward theorems show that first-order theories
with infinite models often possess models of many different sizes.

Therefore first-order language has limited ability to uniquely characterize
infinite structures up to isomorphism.

%%%%%%%%%%%%%%%%%%%%%%%%%%%%%%%%%%%%%%%%%%%%%%%%%%%%%%%%%%%%
\subsection{Skolem Paradox}
\label{subsec:skolem-paradox}
%%%%%%%%%%%%%%%%%%%%%%%%%%%%%%%%%%%%%%%%%%%%%%%%%%%%%%%%%%%%

Set theory proves the existence of uncountable sets such as

\[
\mathbb{R}.
\]

Yet the downward Löwenheim--Skolem theorem implies that, under suitable
metatheoretic assumptions, there can be countable models of substantial
fragments of set theory.

How can a countable model contain what it regards as an uncountable set?

The resolution depends on internal versus external viewpoints.

%%%%%%%%%%%%%%%%%%%%%%%%%%%%%%%%%%%%%%%%%%%%%%%%%%%%%%%%%%%%
\subsection{Internal and External Cardinality}
\label{subsec:internal-external-cardinality}
%%%%%%%%%%%%%%%%%%%%%%%%%%%%%%%%%%%%%%%%%%%%%%%%%%%%%%%%%%%%

Suppose \(\mathcal{M}\) is countable externally.

Inside \(\mathcal{M}\), a set \(R\) may satisfy:

\[
\boxed{
\mathcal{M}
\models
\text{``\(R\) is uncountable.''}
}
\]

This means \(\mathcal{M}\) contains no function that it recognizes as a
bijection

\[
\omega^\mathcal{M}
\to
R.
\]

An external observer may nevertheless enumerate the elements of \(R\).

The external enumeration need not be an element of \(\mathcal{M}\).

%%%%%%%%%%%%%%%%%%%%%%%%%%%%%%%%%%%%%%%%%%%%%%%%%%%%%%%%%%%%
\subsection{Categoricity}
\label{subsec:model-theoretic-categoricity}
%%%%%%%%%%%%%%%%%%%%%%%%%%%%%%%%%%%%%%%%%%%%%%%%%%%%%%%%%%%%

\begin{definitionbox}{Categorical Theory}
A theory \(T\) is categorical in cardinality \(\kappa\) if all models of
\(T\) of cardinality \(\kappa\) are isomorphic.
\end{definitionbox}

Thus if

\[
\mathcal{M},
\mathcal{N}\models T
\]

and

\[
|\mathcal{M}|
=
|\mathcal{N}|
=
\kappa,
\]

then

\[
\boxed{
\mathcal{M}
\cong
\mathcal{N}.
}
\]

%%%%%%%%%%%%%%%%%%%%%%%%%%%%%%%%%%%%%%%%%%%%%%%%%%%%%%%%%%%%
\subsection{Finite Categoricity}
\label{subsec:finite-categoricity}
%%%%%%%%%%%%%%%%%%%%%%%%%%%%%%%%%%%%%%%%%%%%%%%%%%%%%%%%%%%%

Finite structures can often be characterized completely by first-order
sentences.

For example, a specific finite graph can be described up to isomorphism by
a sentence listing:

\[
\boxed{
\text{its finite number of vertices},
}
\]

and

\[
\boxed{
\text{all adjacency relations}.
}
\]

Infinite categoricity behaves very differently.

%%%%%%%%%%%%%%%%%%%%%%%%%%%%%%%%%%%%%%%%%%%%%%%%%%%%%%%%%%%%
\subsection{Infinite Categoricity}
\label{subsec:infinite-categoricity}
%%%%%%%%%%%%%%%%%%%%%%%%%%%%%%%%%%%%%%%%%%%%%%%%%%%%%%%%%%%%

A theory may be categorical in one infinite cardinal while having
nonisomorphic models in another.

This makes cardinal-specific classification important.

Categoricity results are among the central themes of modern model theory.

%%%%%%%%%%%%%%%%%%%%%%%%%%%%%%%%%%%%%%%%%%%%%%%%%%%%%%%%%%%%
\subsection{Complete Theory of Dense Linear Orders}
\label{subsec:dense-linear-orders}
%%%%%%%%%%%%%%%%%%%%%%%%%%%%%%%%%%%%%%%%%%%%%%%%%%%%%%%%%%%%

Consider the language

\[
L=\{<\}.
\]

The theory of dense linear orders without endpoints states, roughly:

\begin{itemize}
    \item \(<\) is a linear order;
    \item between any two distinct elements lies another;
    \item there is no least element;
    \item there is no greatest element.
\end{itemize}

The ordered rationals

\[
\boxed{
(\mathbb{Q},<)
}
\]

are a canonical model.

%%%%%%%%%%%%%%%%%%%%%%%%%%%%%%%%%%%%%%%%%%%%%%%%%%%%%%%%%%%%
\subsection{Countable Categoricity of Dense Linear Orders}
\label{subsec:dlo-countable-categoricity}
%%%%%%%%%%%%%%%%%%%%%%%%%%%%%%%%%%%%%%%%%%%%%%%%%%%%%%%%%%%%

Any two countable dense linear orders without endpoints are isomorphic.

Thus the theory is categorical in

\[
\boxed{
\aleph_0.
}
\]

A back-and-forth argument constructs the isomorphism.

%%%%%%%%%%%%%%%%%%%%%%%%%%%%%%%%%%%%%%%%%%%%%%%%%%%%%%%%%%%%
\subsection{Back-and-Forth Method}
\label{subsec:back-and-forth}
%%%%%%%%%%%%%%%%%%%%%%%%%%%%%%%%%%%%%%%%%%%%%%%%%%%%%%%%%%%%

Suppose

\[
\mathcal{M}
\]

and

\[
\mathcal{N}
\]

are countable structures.

Enumerate:

\[
M
=
\{m_0,m_1,\ldots\},
\]

\[
N
=
\{n_0,n_1,\ldots\}.
\]

Construct increasingly large finite partial isomorphisms.

At alternating stages:

\[
\boxed{
\text{forth: include the next element of }M,
}
\]

and

\[
\boxed{
\text{back: include the next element of }N.
}
\]

The union can yield a full isomorphism.

%%%%%%%%%%%%%%%%%%%%%%%%%%%%%%%%%%%%%%%%%%%%%%%%%%%%%%%%%%%%
\subsection{Worked Example: Back-and-Forth Intuition}
\label{subsec:worked-back-and-forth}
%%%%%%%%%%%%%%%%%%%%%%%%%%%%%%%%%%%%%%%%%%%%%%%%%%%%%%%%%%%%

\begin{workedexamplebox}{Dense Orders}

Suppose finite elements of two countable dense orders have already been
matched order-preservingly.

Take a new point \(a\) in the first order.

Its position is determined relative to finitely many already matched
points.

Because the second order is dense and has no endpoints, choose a new point
\(b\) in the corresponding interval.

This extends the partial isomorphism.

Repeating back and forth produces

\begin{answerbox}
\[
\boxed{
\text{an isomorphism between the two countable dense orders.}
}
\]
\end{answerbox}

\end{workedexamplebox}

%%%%%%%%%%%%%%%%%%%%%%%%%%%%%%%%%%%%%%%%%%%%%%%%%%%%%%%%%%%%
\subsection{Quantifier Elimination}
\label{subsec:quantifier-elimination}
%%%%%%%%%%%%%%%%%%%%%%%%%%%%%%%%%%%%%%%%%%%%%%%%%%%%%%%%%%%%

A theory \(T\) admits quantifier elimination if every formula is equivalent
modulo \(T\) to a quantifier-free formula.

Thus for every

\[
\varphi(\mathbf{x}),
\]

there exists quantifier-free

\[
\psi(\mathbf{x})
\]

such that

\[
\boxed{
T\models
\forall\mathbf{x}
\,
(
\varphi(\mathbf{x})
\leftrightarrow
\psi(\mathbf{x})
).
}
\]

%%%%%%%%%%%%%%%%%%%%%%%%%%%%%%%%%%%%%%%%%%%%%%%%%%%%%%%%%%%%
\subsection{Why Quantifier Elimination Matters}
\label{subsec:why-quantifier-elimination}
%%%%%%%%%%%%%%%%%%%%%%%%%%%%%%%%%%%%%%%%%%%%%%%%%%%%%%%%%%%%

Quantifier elimination can provide:

\[
\boxed{
\text{decidability},
}
\]

\[
\boxed{
\text{explicit definable sets},
}
\]

and

\[
\boxed{
\text{structural understanding of models}.
}
\]

It reduces logical complexity by transforming quantified formulas into
simpler forms.

%%%%%%%%%%%%%%%%%%%%%%%%%%%%%%%%%%%%%%%%%%%%%%%%%%%%%%%%%%%%
\subsection{Example: Dense Linear Orders}
\label{subsec:dlo-quantifier-elimination}
%%%%%%%%%%%%%%%%%%%%%%%%%%%%%%%%%%%%%%%%%%%%%%%%%%%%%%%%%%%%

Dense linear orders without endpoints admit quantifier elimination in the
language of order.

Therefore definable subsets of one variable are built from finite Boolean
combinations of inequalities involving parameters.

This strongly constrains the complexity of definable sets.

%%%%%%%%%%%%%%%%%%%%%%%%%%%%%%%%%%%%%%%%%%%%%%%%%%%%%%%%%%%%
\subsection{Model Completeness}
\label{subsec:model-completeness}
%%%%%%%%%%%%%%%%%%%%%%%%%%%%%%%%%%%%%%%%%%%%%%%%%%%%%%%%%%%%

A theory \(T\) is model complete if every embedding between models of \(T\)
is elementary.

Thus if

\[
\mathcal{M},
\mathcal{N}\models T
\]

and

\[
\mathcal{M}
\subseteq
\mathcal{N},
\]

then

\[
\boxed{
\mathcal{M}
\preccurlyeq
\mathcal{N}
}
\]

whenever the inclusion is an embedding.

Quantifier elimination often implies model completeness.

%%%%%%%%%%%%%%%%%%%%%%%%%%%%%%%%%%%%%%%%%%%%%%%%%%%%%%%%%%%%
\subsection{Prime Models}
\label{subsec:prime-models}
%%%%%%%%%%%%%%%%%%%%%%%%%%%%%%%%%%%%%%%%%%%%%%%%%%%%%%%%%%%%

A model \(\mathcal{M}\models T\) is prime if it elementarily embeds into
every model of \(T\).

Conceptually,

\[
\boxed{
\mathcal{M}
\hookrightarrow_{\mathrm{elem}}
\mathcal{N}
}
\]

for all

\[
\mathcal{N}\models T.
\]

Prime models are minimal in an elementary sense.

%%%%%%%%%%%%%%%%%%%%%%%%%%%%%%%%%%%%%%%%%%%%%%%%%%%%%%%%%%%%
\subsection{Saturated Models}
\label{subsec:saturated-models}
%%%%%%%%%%%%%%%%%%%%%%%%%%%%%%%%%%%%%%%%%%%%%%%%%%%%%%%%%%%%

A saturated model realizes as many consistent types as its size permits.

Roughly, a \(\kappa\)-saturated model realizes every type over fewer than
\(\kappa\) parameters that is consistent with the theory.

Saturated models are rich enough to contain solutions to many locally
consistent descriptions.

%%%%%%%%%%%%%%%%%%%%%%%%%%%%%%%%%%%%%%%%%%%%%%%%%%%%%%%%%%%%
\subsection{Why Saturation Matters}
\label{subsec:why-saturation-matters}
%%%%%%%%%%%%%%%%%%%%%%%%%%%%%%%%%%%%%%%%%%%%%%%%%%%%%%%%%%%%

Saturated models are useful because they support:

\[
\boxed{
\text{uniform type realization},
}
\]

\[
\boxed{
\text{homogeneity},
}
\]

and

\[
\boxed{
\text{classification arguments}.
}
\]

They often behave as universal environments for studying a complete
theory.

%%%%%%%%%%%%%%%%%%%%%%%%%%%%%%%%%%%%%%%%%%%%%%%%%%%%%%%%%%%%
\subsection{Homogeneous Structures}
\label{subsec:homogeneous-structures}
%%%%%%%%%%%%%%%%%%%%%%%%%%%%%%%%%%%%%%%%%%%%%%%%%%%%%%%%%%%%

A structure is homogeneous when small partial symmetries can often be
extended to global automorphisms.

Very roughly,

\[
\boxed{
\text{locally indistinguishable tuples}
}
\]

can be mapped to one another by automorphisms.

Homogeneity is closely related to types and saturation.

%%%%%%%%%%%%%%%%%%%%%%%%%%%%%%%%%%%%%%%%%%%%%%%%%%%%%%%%%%%%
\subsection{Automorphisms}
\label{subsec:model-theoretic-automorphisms}
%%%%%%%%%%%%%%%%%%%%%%%%%%%%%%%%%%%%%%%%%%%%%%%%%%%%%%%%%%%%

An automorphism of \(\mathcal{M}\) is an isomorphism

\[
\boxed{
f:\mathcal{M}\to\mathcal{M}.
}
\]

Automorphisms reveal structural symmetries.

Definable objects without parameters must be preserved by every
automorphism.

This provides an important test for nondefinability.

%%%%%%%%%%%%%%%%%%%%%%%%%%%%%%%%%%%%%%%%%%%%%%%%%%%%%%%%%%%%
\subsection{Worked Example: Automorphism and Definability}
\label{subsec:worked-automorphism-definability}
%%%%%%%%%%%%%%%%%%%%%%%%%%%%%%%%%%%%%%%%%%%%%%%%%%%%%%%%%%%%

\begin{workedexamplebox}{Why Symmetry Restricts Definability}

Suppose structure \(\mathcal{M}\) has an automorphism \(f\) with

\[
f(a)=b,
\qquad
a\neq b.
\]

If \(a\) were uniquely definable without parameters by formula
\(\varphi(x)\), then automorphisms would preserve \(\varphi\).

Thus \(b\) would also satisfy \(\varphi\), contradicting uniqueness.

\begin{answerbox}
\[
\boxed{
a
\text{ cannot be parameter-free uniquely definable if an automorphism
moves it.}
}
\]
\end{answerbox}

\end{workedexamplebox}

%%%%%%%%%%%%%%%%%%%%%%%%%%%%%%%%%%%%%%%%%%%%%%%%%%%%%%%%%%%%
\subsection{Ultrafilters}
\label{subsec:model-theoretic-ultrafilters}
%%%%%%%%%%%%%%%%%%%%%%%%%%%%%%%%%%%%%%%%%%%%%%%%%%%%%%%%%%%%

An ultrafilter \(\mathcal{U}\) on index set \(I\) is a maximal proper
filter.

For every

\[
A\subseteq I,
\]

exactly one of

\[
\boxed{
A\in\mathcal{U}
}
\]

or

\[
\boxed{
I\setminus A\in\mathcal{U}
}
\]

holds.

Ultrafilters allow model theory to convert statements true on
``\(\mathcal{U}\)-many'' indices into statements about a new structure.

%%%%%%%%%%%%%%%%%%%%%%%%%%%%%%%%%%%%%%%%%%%%%%%%%%%%%%%%%%%%
\subsection{Ultraproducts}
\label{subsec:ultraproducts}
%%%%%%%%%%%%%%%%%%%%%%%%%%%%%%%%%%%%%%%%%%%%%%%%%%%%%%%%%%%%

Let

\[
\{
\mathcal{M}_i
\}_{i\in I}
\]

be structures in the same language and let \(\mathcal{U}\) be an
ultrafilter on \(I\).

Start with the Cartesian product

\[
\prod_{i\in I}M_i.
\]

Define two sequences equivalent when

\[
\boxed{
\{
i:
a_i=b_i
\}
\in\mathcal{U}.
}
\]

The quotient structure is the ultraproduct

\[
\boxed{
\prod_{i\in I}
\mathcal{M}_i
/
\mathcal{U}.
}
\]

%%%%%%%%%%%%%%%%%%%%%%%%%%%%%%%%%%%%%%%%%%%%%%%%%%%%%%%%%%%%
\subsection{Ultrapowers}
\label{subsec:ultrapowers}
%%%%%%%%%%%%%%%%%%%%%%%%%%%%%%%%%%%%%%%%%%%%%%%%%%%%%%%%%%%%

If every structure in the ultraproduct is the same structure
\(\mathcal{M}\), then the resulting object is an ultrapower:

\[
\boxed{
\mathcal{M}^I/\mathcal{U}.
}
\]

Ultrapowers produce elementary extensions and are central in nonstandard
analysis and model theory.

%%%%%%%%%%%%%%%%%%%%%%%%%%%%%%%%%%%%%%%%%%%%%%%%%%%%%%%%%%%%
\subsection{Łoś's Theorem}
\label{subsec:los-theorem}
%%%%%%%%%%%%%%%%%%%%%%%%%%%%%%%%%%%%%%%%%%%%%%%%%%%%%%%%%%%%

\begin{theorem}
Let

\[
\mathcal{M}
=
\prod_{i\in I}
\mathcal{M}_i/\mathcal{U}.
\]

For a first-order formula

\[
\varphi(\mathbf{x}),
\]

the ultraproduct satisfies

\[
\varphi([\mathbf{a}_i])
\]

if and only if the set of indices where

\[
\mathcal{M}_i\models\varphi(\mathbf{a}_i)
\]

belongs to \(\mathcal{U}\).
\end{theorem}

Symbolically,

\[
\boxed{
\mathcal{M}
\models
\varphi([\mathbf{a}_i])
\iff
\{
i:
\mathcal{M}_i\models\varphi(\mathbf{a}_i)
\}
\in\mathcal{U}.
}
\]

%%%%%%%%%%%%%%%%%%%%%%%%%%%%%%%%%%%%%%%%%%%%%%%%%%%%%%%%%%%%
\subsection{Why Łoś's Theorem Matters}
\label{subsec:why-los-theorem}
%%%%%%%%%%%%%%%%%%%%%%%%%%%%%%%%%%%%%%%%%%%%%%%%%%%%%%%%%%%%

Łoś's theorem shows that first-order truth transfers through
ultraproducts.

Therefore if every \(\mathcal{M}_i\) satisfies theory \(T\), then the
ultraproduct also satisfies \(T\):

\[
\boxed{
\mathcal{M}_i\models T
\text{ for all }i
\Longrightarrow
\prod_i\mathcal{M}_i/\mathcal{U}
\models T.
}
\]

%%%%%%%%%%%%%%%%%%%%%%%%%%%%%%%%%%%%%%%%%%%%%%%%%%%%%%%%%%%%
\subsection{Nonstandard Analysis Preview}
\label{subsec:nonstandard-analysis-preview-model-theory}
%%%%%%%%%%%%%%%%%%%%%%%%%%%%%%%%%%%%%%%%%%%%%%%%%%%%%%%%%%%%

An ultrapower of the real numbers can produce a larger ordered field

\[
\boxed{
{}^*\mathbb{R}
}
\]

containing infinitesimal and infinitely large elements.

By Łoś's theorem, first-order properties of \(\mathbb{R}\) transfer to the
ultrapower.

This provides one model-theoretic foundation for nonstandard analysis.

%%%%%%%%%%%%%%%%%%%%%%%%%%%%%%%%%%%%%%%%%%%%%%%%%%%%%%%%%%%%
\subsection{Infinitesimals}
\label{subsec:model-theoretic-infinitesimals}
%%%%%%%%%%%%%%%%%%%%%%%%%%%%%%%%%%%%%%%%%%%%%%%%%%%%%%%%%%%%

A positive infinitesimal \(\varepsilon\) satisfies

\[
\boxed{
0<\varepsilon<\frac1n
}
\]

for every positive standard natural number \(n\).

Such elements do not exist in the ordinary real field.

They can appear in suitable nonstandard elementary extensions.

%%%%%%%%%%%%%%%%%%%%%%%%%%%%%%%%%%%%%%%%%%%%%%%%%%%%%%%%%%%%
\subsection{Theory of Algebraically Closed Fields}
\label{subsec:algebraically-closed-fields-model-theory}
%%%%%%%%%%%%%%%%%%%%%%%%%%%%%%%%%%%%%%%%%%%%%%%%%%%%%%%%%%%%

Model theory studies algebraically closed fields through first-order
properties.

For a fixed characteristic, the theory of algebraically closed fields is a
central example.

It possesses strong model-theoretic properties including:

\[
\boxed{
\text{completeness},
}
\]

and

\[
\boxed{
\text{quantifier elimination}
}
\]

in the standard field language.

This connects model theory directly with algebraic geometry.

%%%%%%%%%%%%%%%%%%%%%%%%%%%%%%%%%%%%%%%%%%%%%%%%%%%%%%%%%%%%
\subsection{Definable Sets in Algebraically Closed Fields}
\label{subsec:definable-sets-acf}
%%%%%%%%%%%%%%%%%%%%%%%%%%%%%%%%%%%%%%%%%%%%%%%%%%%%%%%%%%%%

Quantifier elimination implies that definable sets can be described by
Boolean combinations of polynomial equations.

Thus the model-theoretic notion of definability is closely connected to
constructible sets in algebraic geometry.

This is one of the foundational bridges between logic and geometry.

%%%%%%%%%%%%%%%%%%%%%%%%%%%%%%%%%%%%%%%%%%%%%%%%%%%%%%%%%%%%
\subsection{Real Closed Fields}
\label{subsec:real-closed-fields-model-theory}
%%%%%%%%%%%%%%%%%%%%%%%%%%%%%%%%%%%%%%%%%%%%%%%%%%%%%%%%%%%%

The real numbers form a real closed field.

The first-order theory of real closed fields has strong properties,
including quantifier elimination in the language of ordered fields.

This implies decidability of first-order statements about ordered-field
operations.

%%%%%%%%%%%%%%%%%%%%%%%%%%%%%%%%%%%%%%%%%%%%%%%%%%%%%%%%%%%%
\subsection{Semialgebraic Sets}
\label{subsec:semialgebraic-model-theory}
%%%%%%%%%%%%%%%%%%%%%%%%%%%%%%%%%%%%%%%%%%%%%%%%%%%%%%%%%%%%

A semialgebraic subset of

\[
\mathbb{R}^n
\]

is constructed from polynomial equalities and inequalities using finite
Boolean operations.

Examples include

\[
\boxed{
x^2+y^2\leq1,
}
\]

and

\[
\boxed{
x>0
\land
y^2<x.
}
\]

Quantifier elimination for real closed fields shows that projections of
semialgebraic sets remain semialgebraic.

%%%%%%%%%%%%%%%%%%%%%%%%%%%%%%%%%%%%%%%%%%%%%%%%%%%%%%%%%%%%
\subsection{Tarski--Seidenberg Principle}
\label{subsec:tarski-seidenberg}
%%%%%%%%%%%%%%%%%%%%%%%%%%%%%%%%%%%%%%%%%%%%%%%%%%%%%%%%%%%%

A central consequence of real quantifier elimination is that the projection
of a semialgebraic set is semialgebraic.

If

\[
S
\subseteq
\mathbb{R}^{n+m}
\]

is semialgebraic, then its projection onto

\[
\mathbb{R}^n
\]

is also semialgebraic.

This has applications in:

\[
\boxed{
\text{real algebraic geometry},
}
\]

\[
\boxed{
\text{optimization},
}
\]

and

\[
\boxed{
\text{control theory}.
}
\]

%%%%%%%%%%%%%%%%%%%%%%%%%%%%%%%%%%%%%%%%%%%%%%%%%%%%%%%%%%%%
\subsection{O-Minimal Structures Preview}
\label{subsec:o-minimality-preview}
%%%%%%%%%%%%%%%%%%%%%%%%%%%%%%%%%%%%%%%%%%%%%%%%%%%%%%%%%%%%

An ordered structure is o-minimal if every definable subset of its
one-dimensional domain is a finite union of points and intervals.

Thus definable one-dimensional sets are geometrically tame.

O-minimality provides a framework for studying structures that expand the
real field while retaining controlled geometry.

%%%%%%%%%%%%%%%%%%%%%%%%%%%%%%%%%%%%%%%%%%%%%%%%%%%%%%%%%%%%
\subsection{Why O-Minimality Matters}
\label{subsec:why-o-minimality-matters}
%%%%%%%%%%%%%%%%%%%%%%%%%%%%%%%%%%%%%%%%%%%%%%%%%%%%%%%%%%%%

O-minimal structures rule out highly oscillatory or fractal definable
subsets of the line.

They support strong results concerning:

\[
\boxed{
\text{dimension},
}
\]

\[
\boxed{
\text{cell decomposition},
}
\]

\[
\boxed{
\text{differentiability},
}
\]

and

\[
\boxed{
\text{tame optimization}.
}
\]

%%%%%%%%%%%%%%%%%%%%%%%%%%%%%%%%%%%%%%%%%%%%%%%%%%%%%%%%%%%%
\subsection{Model Theory of Graphs}
\label{subsec:model-theory-graphs}
%%%%%%%%%%%%%%%%%%%%%%%%%%%%%%%%%%%%%%%%%%%%%%%%%%%%%%%%%%%%

Graphs can be treated as structures

\[
\boxed{
\mathcal{G}
=
(V,E),
}
\]

where \(E\) is a binary relation.

First-order graph formulas can express properties such as:

\[
\boxed{
\text{adjacency},
}
\]

\[
\boxed{
\text{existence of triangles},
}
\]

and

\[
\boxed{
\text{existence of vertices with specified local configurations}.
}
\]

%%%%%%%%%%%%%%%%%%%%%%%%%%%%%%%%%%%%%%%%%%%%%%%%%%%%%%%%%%%%
\subsection{Finite Model Theory}
\label{subsec:finite-model-theory}
%%%%%%%%%%%%%%%%%%%%%%%%%%%%%%%%%%%%%%%%%%%%%%%%%%%%%%%%%%%%

Finite model theory restricts attention to finite structures.

Many classical first-order theorems behave differently in this setting.

For example, compactness does not hold when only finite models are allowed.

Finite model theory is closely connected with:

\[
\boxed{
\text{databases},
}
\]

\[
\boxed{
\text{descriptive complexity},
}
\]

and

\[
\boxed{
\text{theoretical computer science}.
}
\]

%%%%%%%%%%%%%%%%%%%%%%%%%%%%%%%%%%%%%%%%%%%%%%%%%%%%%%%%%%%%
\subsection{Failure of Finite Compactness}
\label{subsec:finite-compactness-failure}
%%%%%%%%%%%%%%%%%%%%%%%%%%%%%%%%%%%%%%%%%%%%%%%%%%%%%%%%%%%%

Consider a theory containing sentences

\[
\sigma_n
\]

asserting that there are at least \(n\) elements for every \(n\).

Every finite subset has a finite model.

But the full theory has no finite model.

Thus compactness fails if the word ``model'' is restricted to finite
structures.

%%%%%%%%%%%%%%%%%%%%%%%%%%%%%%%%%%%%%%%%%%%%%%%%%%%%%%%%%%%%
\subsection{Stability Theory Preview}
\label{subsec:stability-theory-preview}
%%%%%%%%%%%%%%%%%%%%%%%%%%%%%%%%%%%%%%%%%%%%%%%%%%%%%%%%%%%%

Stability theory studies how complicated the space of types of a complete
theory can be.

Roughly, stable theories have controlled numbers of types over parameter
sets.

This provides a classification framework for first-order theories.

%%%%%%%%%%%%%%%%%%%%%%%%%%%%%%%%%%%%%%%%%%%%%%%%%%%%%%%%%%%%
\subsection{The Order Property}
\label{subsec:model-theoretic-order-property}
%%%%%%%%%%%%%%%%%%%%%%%%%%%%%%%%%%%%%%%%%%%%%%%%%%%%%%%%%%%%

A theory exhibits the order property when a formula can encode arbitrarily
long finite linear orders in a uniform way.

Such behavior indicates instability.

Stable theories avoid this form of combinatorial complexity.

%%%%%%%%%%%%%%%%%%%%%%%%%%%%%%%%%%%%%%%%%%%%%%%%%%%%%%%%%%%%
\subsection{Classification Theory}
\label{subsec:model-theoretic-classification-theory}
%%%%%%%%%%%%%%%%%%%%%%%%%%%%%%%%%%%%%%%%%%%%%%%%%%%%%%%%%%%%

Classification theory asks:

\[
\boxed{
\text{Which theories have well-behaved model classes?}
}
\]

Important dividing lines include notions related to:

\[
\boxed{
\text{stability},
}
\]

\[
\boxed{
\text{simplicity},
}
\]

and

\[
\boxed{
\text{dependence}.
}
\]

The goal is not merely to count models but to understand structural
complexity.

%%%%%%%%%%%%%%%%%%%%%%%%%%%%%%%%%%%%%%%%%%%%%%%%%%%%%%%%%%%%
\subsection{Model-Theoretic Independence}
\label{subsec:model-theoretic-independence}
%%%%%%%%%%%%%%%%%%%%%%%%%%%%%%%%%%%%%%%%%%%%%%%%%%%%%%%%%%%%

In stable theories, algebraic dependence can be generalized through
model-theoretic independence relations.

These relations behave analogously to:

\[
\boxed{
\text{linear independence},
}
\]

or

\[
\boxed{
\text{algebraic independence}.
}
\]

They provide a geometry of types and parameter dependence.

%%%%%%%%%%%%%%%%%%%%%%%%%%%%%%%%%%%%%%%%%%%%%%%%%%%%%%%%%%%%
\subsection{Types as Generalized Coordinates}
\label{subsec:types-generalized-coordinates}
%%%%%%%%%%%%%%%%%%%%%%%%%%%%%%%%%%%%%%%%%%%%%%%%%%%%%%%%%%%%

A type records every first-order property of an element relative to chosen
parameters.

In this sense, types act like logical coordinates.

Two tuples with the same type over parameter set \(A\) cannot be
distinguished by formulas using parameters from \(A\).

%%%%%%%%%%%%%%%%%%%%%%%%%%%%%%%%%%%%%%%%%%%%%%%%%%%%%%%%%%%%
\subsection{Automorphisms and Types}
\label{subsec:automorphisms-and-types}
%%%%%%%%%%%%%%%%%%%%%%%%%%%%%%%%%%%%%%%%%%%%%%%%%%%%%%%%%%%%

In sufficiently homogeneous models, tuples having the same type over \(A\)
can often be mapped to one another by an automorphism fixing \(A\).

Thus

\[
\boxed{
\text{same type}
}
\]

is closely connected to

\[
\boxed{
\text{structural symmetry}.
}
\]

%%%%%%%%%%%%%%%%%%%%%%%%%%%%%%%%%%%%%%%%%%%%%%%%%%%%%%%%%%%%
\subsection{Fraïssé Theory Preview}
\label{subsec:fraisse-theory-preview}
%%%%%%%%%%%%%%%%%%%%%%%%%%%%%%%%%%%%%%%%%%%%%%%%%%%%%%%%%%%%

Fraïssé theory constructs countable homogeneous structures from classes of
finite structures.

A suitable class satisfies properties involving:

\[
\boxed{
\text{substructures},
}
\]

\[
\boxed{
\text{joint embeddings},
}
\]

and

\[
\boxed{
\text{amalgamation}.
}
\]

Its limit is a countable structure containing all finite structures in the
class in a highly homogeneous way.

%%%%%%%%%%%%%%%%%%%%%%%%%%%%%%%%%%%%%%%%%%%%%%%%%%%%%%%%%%%%
\subsection{Amalgamation Property}
\label{subsec:amalgamation-property}
%%%%%%%%%%%%%%%%%%%%%%%%%%%%%%%%%%%%%%%%%%%%%%%%%%%%%%%%%%%%

Suppose structure \(A\) embeds into \(B\) and \(C\).

The amalgamation property asks whether \(B\) and \(C\) can be embedded into
a larger structure \(D\) so that the two copies of \(A\) agree.

Conceptually,

\[
\boxed{
B
\leftarrow
A
\rightarrow
C
}
\]

can be completed to

\[
\boxed{
B
\rightarrow
D
\leftarrow
C.
}
\]

%%%%%%%%%%%%%%%%%%%%%%%%%%%%%%%%%%%%%%%%%%%%%%%%%%%%%%%%%%%%
\subsection{Random Graph Preview}
\label{subsec:random-graph-model-theory}
%%%%%%%%%%%%%%%%%%%%%%%%%%%%%%%%%%%%%%%%%%%%%%%%%%%%%%%%%%%%

A famous countable homogeneous graph can be characterized by an extension
property:

for any finite disjoint sets of vertices \(A\) and \(B\), there exists a
vertex adjacent to every vertex in \(A\) and to no vertex in \(B\).

This property determines the countable graph up to isomorphism.

It is a central example of back-and-forth and Fraïssé-style reasoning.

%%%%%%%%%%%%%%%%%%%%%%%%%%%%%%%%%%%%%%%%%%%%%%%%%%%%%%%%%%%%
\subsection{Model Theory and Algebra}
\label{subsec:model-theory-algebra}
%%%%%%%%%%%%%%%%%%%%%%%%%%%%%%%%%%%%%%%%%%%%%%%%%%%%%%%%%%%%

Model-theoretic methods can study:

\[
\boxed{
\text{fields},
}
\]

\[
\boxed{
\text{groups},
}
\]

\[
\boxed{
\text{modules},
}
\]

and

\[
\boxed{
\text{rings}.
}
\]

Questions include:

\[
\boxed{
\text{Which subsets are definable?}
}
\]

\[
\boxed{
\text{Which theories are complete?}
}
\]

and

\[
\boxed{
\text{How many nonisomorphic models exist?}
}
\]

%%%%%%%%%%%%%%%%%%%%%%%%%%%%%%%%%%%%%%%%%%%%%%%%%%%%%%%%%%%%
\subsection{Model Theory and Number Theory}
\label{subsec:model-theory-number-theory}
%%%%%%%%%%%%%%%%%%%%%%%%%%%%%%%%%%%%%%%%%%%%%%%%%%%%%%%%%%%%

Interactions with number theory include model-theoretic study of:

\[
\boxed{
\text{fields},
}
\]

\[
\boxed{
\text{valuations},
}
\]

\[
\boxed{
\text{Diophantine equations},
}
\]

and

\[
\boxed{
\text{definability in arithmetic structures}.
}
\]

Logical properties can reveal arithmetic information and vice versa.

%%%%%%%%%%%%%%%%%%%%%%%%%%%%%%%%%%%%%%%%%%%%%%%%%%%%%%%%%%%%
\subsection{Model Theory and Geometry}
\label{subsec:model-theory-geometry}
%%%%%%%%%%%%%%%%%%%%%%%%%%%%%%%%%%%%%%%%%%%%%%%%%%%%%%%%%%%%

Definable sets often possess geometric structure.

In algebraically closed fields, definability is linked to algebraic
geometry.

In real closed and o-minimal structures, definability is linked to tame
real geometry.

Thus model theory can serve as a language for generalized geometry.

%%%%%%%%%%%%%%%%%%%%%%%%%%%%%%%%%%%%%%%%%%%%%%%%%%%%%%%%%%%%
\subsection{Model Theory and Analysis}
\label{subsec:model-theory-analysis}
%%%%%%%%%%%%%%%%%%%%%%%%%%%%%%%%%%%%%%%%%%%%%%%%%%%%%%%%%%%%

Extensions of classical model theory study metric and analytic structures.

Applications include structures arising from:

\[
\boxed{
\text{Banach spaces},
}
\]

\[
\boxed{
\text{probability spaces},
}
\]

and

\[
\boxed{
\text{operator algebras}.
}
\]

These settings often require variants such as continuous logic.

%%%%%%%%%%%%%%%%%%%%%%%%%%%%%%%%%%%%%%%%%%%%%%%%%%%%%%%%%%%%
\subsection{Model Theory and Databases}
\label{subsec:model-theory-databases}
%%%%%%%%%%%%%%%%%%%%%%%%%%%%%%%%%%%%%%%%%%%%%%%%%%%%%%%%%%%%

A finite relational database can be viewed as a finite structure.

Database queries correspond naturally to logical formulas.

Thus questions about query expressiveness become questions about definability
in finite structures.

This connects finite model theory with database theory.

%%%%%%%%%%%%%%%%%%%%%%%%%%%%%%%%%%%%%%%%%%%%%%%%%%%%%%%%%%%%
\subsection{Model Theory and Computation}
\label{subsec:model-theory-computation}
%%%%%%%%%%%%%%%%%%%%%%%%%%%%%%%%%%%%%%%%%%%%%%%%%%%%%%%%%%%%

Model theory interacts with computation through:

\[
\boxed{
\text{decidable theories},
}
\]

\[
\boxed{
\text{quantifier elimination},
}
\]

\[
\boxed{
\text{finite structures},
}
\]

and

\[
\boxed{
\text{automated reasoning}.
}
\]

A structural theorem about models can sometimes lead directly to an
algorithm.

%%%%%%%%%%%%%%%%%%%%%%%%%%%%%%%%%%%%%%%%%%%%%%%%%%%%%%%%%%%%
\subsection{Model-Theoretic Workflow}
\label{subsec:model-theory-workflow}
%%%%%%%%%%%%%%%%%%%%%%%%%%%%%%%%%%%%%%%%%%%%%%%%%%%%%%%%%%%%

When studying a theory model-theoretically:

\begin{enumerate}
    \item specify the first-order language;
    \item identify the axioms of the theory;
    \item characterize important models;
    \item determine whether the theory has finite or infinite models;
    \item compare structures up to isomorphism;
    \item compare structures up to elementary equivalence;
    \item identify embeddings and substructures;
    \item test whether substructures are elementary;
    \item describe definable sets and functions;
    \item identify definable and algebraic closure;
    \item determine whether the theory is complete;
    \item investigate compactness consequences;
    \item apply Löwenheim--Skolem where relevant;
    \item analyze model cardinalities;
    \item investigate categoricity;
    \item study quantifier elimination;
    \item study types and their realizations;
    \item construct elementary extensions when useful;
    \item use ultraproducts for new models;
    \item investigate automorphisms and homogeneity;
    \item classify model complexity using stability-style ideas when
          appropriate;
    \item distinguish internal from external statements about models.
\end{enumerate}

%%%%%%%%%%%%%%%%%%%%%%%%%%%%%%%%%%%%%%%%%%%%%%%%%%%%%%%%%%%%
\subsection{Common Errors}
\label{subsec:model-theory-common-errors}
%%%%%%%%%%%%%%%%%%%%%%%%%%%%%%%%%%%%%%%%%%%%%%%%%%%%%%%%%%%%

\subsubsection{Confusing a Theory With One of Its Models}

A theory is a set of sentences.

A model is a structure satisfying those sentences.

Thus

\[
\boxed{
T
\neq
\mathcal{M}.
}
\]

%%%%%%%%%%%%%%%%%%%%%%%%%%%%%%%%%%%%%%%%%%%%%%%%%%%%%%%%%%%%
\subsubsection{Confusing Isomorphism With Elementary Equivalence}

Isomorphic structures are elementarily equivalent.

But elementarily equivalent infinite structures need not be isomorphic.

%%%%%%%%%%%%%%%%%%%%%%%%%%%%%%%%%%%%%%%%%%%%%%%%%%%%%%%%%%%%
\subsubsection{Assuming Every Substructure Is Elementary}

Substructures preserve atomic structure but may fail to preserve quantified
formulas.

Elementary substructure is a stronger condition.

%%%%%%%%%%%%%%%%%%%%%%%%%%%%%%%%%%%%%%%%%%%%%%%%%%%%%%%%%%%%
\subsubsection{Confusing \(\subseteq\) and \(\preccurlyeq\)}

The notation

\[
\mathcal{M}\subseteq\mathcal{N}
\]

means substructure.

The notation

\[
\boxed{
\mathcal{M}\preccurlyeq\mathcal{N}
}
\]

means elementary substructure.

The second implies much more.

%%%%%%%%%%%%%%%%%%%%%%%%%%%%%%%%%%%%%%%%%%%%%%%%%%%%%%%%%%%%
\subsubsection{Ignoring Parameters in Definability}

A set may fail to be definable without parameters but become definable once
specific elements are named.

Always specify whether parameters are allowed.

%%%%%%%%%%%%%%%%%%%%%%%%%%%%%%%%%%%%%%%%%%%%%%%%%%%%%%%%%%%%
\subsubsection{Confusing Logical Completeness With Complete Theory}

Logical completeness concerns proof systems.

Theory completeness concerns whether every sentence is decided.

These are separate notions.

%%%%%%%%%%%%%%%%%%%%%%%%%%%%%%%%%%%%%%%%%%%%%%%%%%%%%%%%%%%%
\subsubsection{Assuming a Finitely Satisfiable Type Is Realized in the Same Model}

Finite satisfiability guarantees consistency with the theory under suitable
conditions.

It does not mean the original model realizes the type.

An elementary extension may be required.

%%%%%%%%%%%%%%%%%%%%%%%%%%%%%%%%%%%%%%%%%%%%%%%%%%%%%%%%%%%%
\subsubsection{Misreading Compactness}

Compactness does not say that an infinite model is finite.

It says that satisfiability of a first-order theory is determined by all
finite subsets of its sentences.

%%%%%%%%%%%%%%%%%%%%%%%%%%%%%%%%%%%%%%%%%%%%%%%%%%%%%%%%%%%%
\subsubsection{Assuming Compactness Applies Only to Finite Structures}

The theorem ranges over ordinary first-order structures.

In fact, restricting attention only to finite structures destroys
compactness.

%%%%%%%%%%%%%%%%%%%%%%%%%%%%%%%%%%%%%%%%%%%%%%%%%%%%%%%%%%%%
\subsubsection{Thinking Nonstandard Models Violate the Axioms}

A nonstandard model satisfies the same first-order axioms.

It differs from the intended structure because those axioms do not
categorically determine the model.

%%%%%%%%%%%%%%%%%%%%%%%%%%%%%%%%%%%%%%%%%%%%%%%%%%%%%%%%%%%%
\subsubsection{Treating Skolem's Paradox as a Contradiction}

A countable model can internally regard a set as uncountable because the
required external enumeration is not an element of the model.

The internal and external viewpoints differ.

%%%%%%%%%%%%%%%%%%%%%%%%%%%%%%%%%%%%%%%%%%%%%%%%%%%%%%%%%%%%
\subsubsection{Assuming First-Order Theories Determine Infinite Cardinality Uniquely}

Löwenheim--Skolem shows that infinite first-order theories often have models
in many cardinalities.

%%%%%%%%%%%%%%%%%%%%%%%%%%%%%%%%%%%%%%%%%%%%%%%%%%%%%%%%%%%%
\subsubsection{Confusing Categoricity With Completeness}

A theory may be complete without being categorical in every cardinal.

Categoricity concerns uniqueness up to isomorphism at a specified
cardinality.

%%%%%%%%%%%%%%%%%%%%%%%%%%%%%%%%%%%%%%%%%%%%%%%%%%%%%%%%%%%%
\subsubsection{Assuming Quantifier Elimination Means Quantifiers Are Removed Syntactically Without a Theory}

Quantifier elimination means equivalence modulo a particular theory.

The transformed quantifier-free formula may depend on the axioms of that
theory.

%%%%%%%%%%%%%%%%%%%%%%%%%%%%%%%%%%%%%%%%%%%%%%%%%%%%%%%%%%%%
\subsubsection{Confusing Algebraic Closure in Fields With Model-Theoretic Algebraic Closure in Every Theory}

The model-theoretic concept generalizes the field-theoretic one.

Its behavior depends on the underlying theory.

%%%%%%%%%%%%%%%%%%%%%%%%%%%%%%%%%%%%%%%%%%%%%%%%%%%%%%%%%%%%
\subsubsection{Assuming Every Ultrafilter Is Principal}

Infinite sets can admit nonprincipal ultrafilters under suitable choice
principles.

These are especially important in ultrapower constructions.

%%%%%%%%%%%%%%%%%%%%%%%%%%%%%%%%%%%%%%%%%%%%%%%%%%%%%%%%%%%%
\subsubsection{Assuming Ultraproduct Equality Is Pointwise Equality}

Two sequences represent the same ultraproduct element when they agree on a
set belonging to the ultrafilter.

They need not agree at every index.

%%%%%%%%%%%%%%%%%%%%%%%%%%%%%%%%%%%%%%%%%%%%%%%%%%%%%%%%%%%%
\subsubsection{Applying Łoś's Theorem to Arbitrary Higher-Order Statements}

Łoś's theorem is fundamentally a first-order transfer theorem.

One must verify that the property being transferred is expressible in the
relevant first-order language.

%%%%%%%%%%%%%%%%%%%%%%%%%%%%%%%%%%%%%%%%%%%%%%%%%%%%%%%%%%%%
\subsubsection{Assuming Stability Means Numerical Stability}

Model-theoretic stability is a classification-theoretic property of a
first-order theory.

It has nothing to do with numerical stability from numerical analysis.

%%%%%%%%%%%%%%%%%%%%%%%%%%%%%%%%%%%%%%%%%%%%%%%%%%%%%%%%%%%%
\subsection{Applications}
\label{subsec:model-theory-applications}
%%%%%%%%%%%%%%%%%%%%%%%%%%%%%%%%%%%%%%%%%%%%%%%%%%%%%%%%%%%%

\begin{applicationbox}

\textbf{Algebra.}
Model theory studies definability, elementary equivalence, and classification
of fields, groups, modules, and rings.

\medskip

\textbf{Algebraic Geometry.}
Quantifier elimination for algebraically closed fields links definable sets
with algebraic and constructible geometry.

\medskip

\textbf{Real Algebraic Geometry.}
Real closed fields and the Tarski--Seidenberg principle connect model theory
with polynomial inequalities and semialgebraic sets.

\medskip

\textbf{Number Theory.}
Definability, valued fields, Diophantine problems, and logical properties of
arithmetic structures form major areas of interaction.

\medskip

\textbf{Analysis.}
Continuous and metric variants of model theory study Banach spaces,
probability structures, and operator-algebraic systems.

\medskip

\textbf{Computer Science.}
Finite model theory, database logic, decidable theories, and automated
reasoning use model-theoretic ideas.

\medskip

\textbf{Foundations.}
Compactness, Löwenheim--Skolem, nonstandard models, ultraproducts, and
categoricity reveal the limitations and possibilities of first-order
axiomatization.

\end{applicationbox}

%%%%%%%%%%%%%%%%%%%%%%%%%%%%%%%%%%%%%%%%%%%%%%%%%%%%%%%%%%%%
\subsection{Exercises}
\label{subsec:model-theory-exercises}
%%%%%%%%%%%%%%%%%%%%%%%%%%%%%%%%%%%%%%%%%%%%%%%%%%%%%%%%%%%%

\begin{exercise}[1]
Define an \(L\)-structure.
\end{exercise}

\begin{exercise}[1]
Define a model of a theory.
\end{exercise}

\begin{exercise}[1]
Define

\[
\operatorname{Mod}(T).
\]
\end{exercise}

\begin{exercise}[1]
Define

\[
\operatorname{Th}(\mathcal{M}).
\]
\end{exercise}

\begin{exercise}[1]
Define elementary equivalence.
\end{exercise}

\begin{exercise}[1]
Define an isomorphism of structures.
\end{exercise}

\begin{exercise}[1]
Define an elementary substructure.
\end{exercise}

\begin{exercise}[1]
Define a definable set.
\end{exercise}

\begin{exercise}[1]
Define a complete theory.
\end{exercise}

\begin{exercise}[1]
Define a type.
\end{exercise}

\begin{exercise}[2]
Give two isomorphic structures in the same language.
\end{exercise}

\begin{exercise}[2]
Prove that isomorphic structures are elementarily equivalent.
\end{exercise}

\begin{exercise}[3]
Give an example illustrating why elementary equivalence need not imply
isomorphism.
\end{exercise}

\begin{exercise}[2]
Define a substructure for a language containing one binary function.
\end{exercise}

\begin{exercise}[3]
Explain why closure under function symbols is required for a substructure.
\end{exercise}

\begin{exercise}[3]
Give an example of a substructure that is not elementary.
\end{exercise}

\begin{exercise}[3]
State the Tarski--Vaught criterion.
\end{exercise}

\begin{exercise}[3]
Explain the witness condition in the Tarski--Vaught criterion.
\end{exercise}

\begin{exercise}[2]
Give a formula defining the positive elements of an ordered field.
\end{exercise}

\begin{exercise}[2]
Give a formula defining the interval

\[
(a,b)
\]

using parameters \(a,b\).
\end{exercise}

\begin{exercise}[3]
Distinguish parameter-free and parameter definability.
\end{exercise}

\begin{exercise}[3]
Explain why the graph of a definable function must be definable.
\end{exercise}

\begin{exercise}[3]
Explain definable closure conceptually.
\end{exercise}

\begin{exercise}[3]
Explain model-theoretic algebraic closure conceptually.
\end{exercise}

\begin{exercise}[3]
Show that

\[
\operatorname{dcl}(A)
\subseteq
\operatorname{acl}(A).
\]
\end{exercise}

\begin{exercise}[2]
Explain the difference between a complete theory and a complete proof
system.
\end{exercise}

\begin{exercise}[3]
Show that

\[
\operatorname{Th}(\mathcal{M})
\]

is complete.
\end{exercise}

\begin{exercise}[2]
Construct a partial type over the ordered real field.
\end{exercise}

\begin{exercise}[3]
Explain what it means for an element to realize a type.
\end{exercise}

\begin{exercise}[3]
Give a type finitely satisfiable in \(\mathbb{N}\) but not realized there.
\end{exercise}

\begin{exercise}[2]
State the Compactness Theorem.
\end{exercise}

\begin{exercise}[3]
State the semantic consequence form of compactness.
\end{exercise}

\begin{exercise}[3]
Use compactness to show that a theory with arbitrarily large finite models
has an infinite model.
\end{exercise}

\begin{exercise}[3]
Use compactness to produce an element greater than every standard numeral
in a model of arithmetic.
\end{exercise}

\begin{exercise}[3]
Explain why the resulting model is nonstandard.
\end{exercise}

\begin{exercise}[3]
Define an elementary extension.
\end{exercise}

\begin{exercise}[3]
Explain the purpose of an elementary diagram.
\end{exercise}

\begin{exercise}[2]
State the downward Löwenheim--Skolem theorem conceptually.
\end{exercise}

\begin{exercise}[2]
State the upward Löwenheim--Skolem theorem conceptually.
\end{exercise}

\begin{exercise}[3]
Explain why Löwenheim--Skolem limits first-order categoricity.
\end{exercise}

\begin{exercise}[3]
Explain the Skolem paradox.
\end{exercise}

\begin{exercise}[3]
Distinguish internal and external countability.
\end{exercise}

\begin{exercise}[2]
Define categoricity in cardinality \(\kappa\).
\end{exercise}

\begin{exercise}[3]
Explain why a finite structure can often be characterized up to
isomorphism.
\end{exercise}

\begin{exercise}[3]
State the countable categoricity of dense linear orders without endpoints.
\end{exercise}

\begin{exercise}[3]
Explain the back-and-forth method.
\end{exercise}

\begin{exercise}[3]
Apply one ``forth'' step to two dense linear orders.
\end{exercise}

\begin{exercise}[3]
Apply one ``back'' step to two dense linear orders.
\end{exercise}

\begin{exercise}[2]
Define quantifier elimination.
\end{exercise}

\begin{exercise}[3]
Explain why quantifier elimination can imply decidability.
\end{exercise}

\begin{exercise}[3]
Explain why quantifier elimination constrains definable sets.
\end{exercise}

\begin{exercise}[3]
Define model completeness.
\end{exercise}

\begin{exercise}[3]
Explain the relationship between quantifier elimination and model
completeness.
\end{exercise}

\begin{exercise}[3]
Define a prime model conceptually.
\end{exercise}

\begin{exercise}[3]
Define saturation conceptually.
\end{exercise}

\begin{exercise}[3]
Explain why saturated models realize many consistent types.
\end{exercise}

\begin{exercise}[3]
Explain homogeneity.
\end{exercise}

\begin{exercise}[2]
Define an automorphism.
\end{exercise}

\begin{exercise}[3]
Explain why parameter-free definable elements are fixed by all
automorphisms.
\end{exercise}

\begin{exercise}[3]
Use an automorphism to show that a specified element is not uniquely
definable.
\end{exercise}

\begin{exercise}[2]
Define an ultrafilter.
\end{exercise}

\begin{exercise}[3]
Describe the equivalence relation used in an ultraproduct.
\end{exercise}

\begin{exercise}[3]
Define an ultrapower.
\end{exercise}

\begin{exercise}[3]
State Łoś's theorem conceptually.
\end{exercise}

\begin{exercise}[3]
Explain why ultraproducts preserve first-order theories.
\end{exercise}

\begin{exercise}[3]
Explain how ultrapowers lead to nonstandard models.
\end{exercise}

\begin{exercise}[3]
Define a positive infinitesimal in a nonstandard ordered field.
\end{exercise}

\begin{exercise}[3]
Explain why no positive infinitesimal exists in \(\mathbb{R}\).
\end{exercise}

\begin{exercise}[3]
Explain why algebraically closed fields are important in model theory.
\end{exercise}

\begin{exercise}[3]
Explain the role of quantifier elimination for algebraically closed fields.
\end{exercise}

\begin{exercise}[3]
Define a semialgebraic set.
\end{exercise}

\begin{exercise}[3]
State the Tarski--Seidenberg principle conceptually.
\end{exercise}

\begin{exercise}[3]
Explain why projections correspond to existential quantifiers.
\end{exercise}

\begin{exercise}[3]
Define o-minimality conceptually.
\end{exercise}

\begin{exercise}[3]
Give three examples of one-dimensional sets allowed in an o-minimal
structure.
\end{exercise}

\begin{exercise}[3]
Explain finite model theory.
\end{exercise}

\begin{exercise}[3]
Show why ordinary compactness fails when restricted to finite models.
\end{exercise}

\begin{exercise}[3]
Explain stability theory conceptually.
\end{exercise}

\begin{exercise}[3]
Explain the order property conceptually.
\end{exercise}

\begin{exercise}[3]
Explain why model-theoretic stability is unrelated to numerical stability.
\end{exercise}

\begin{exercise}[3]
Describe the goal of classification theory.
\end{exercise}

\begin{exercise}[3]
Explain model-theoretic independence conceptually.
\end{exercise}

\begin{exercise}[3]
Explain why types can be viewed as logical coordinates.
\end{exercise}

\begin{exercise}[3]
Explain the connection between types and automorphisms.
\end{exercise}

\begin{exercise}[3]
Describe the amalgamation property.
\end{exercise}

\begin{exercise}[3]
Explain the idea of a Fraïssé limit.
\end{exercise}

\begin{exercise}[3]
State the extension property of the countable random graph.
\end{exercise}

\begin{exercise}[4]
Prove that isomorphic structures satisfy the same first-order formulas.
\end{exercise}

\begin{exercise}[4]
Prove the Tarski--Vaught criterion.
\end{exercise}

\begin{exercise}[4]
Study elementary chains and their unions.
\end{exercise}

\begin{exercise}[4]
Prove the Compactness Theorem from first-order completeness.
\end{exercise}

\begin{exercise}[4]
Derive the downward Löwenheim--Skolem theorem using Skolem functions.
\end{exercise}

\begin{exercise}[4]
Study the upward Löwenheim--Skolem theorem.
\end{exercise}

\begin{exercise}[4]
Construct a nonstandard model of arithmetic using compactness.
\end{exercise}

\begin{exercise}[4]
Study complete types and Stone spaces.
\end{exercise}

\begin{exercise}[4]
Study omitting-types theorems.
\end{exercise}

\begin{exercise}[4]
Study saturated and homogeneous models.
\end{exercise}

\begin{exercise}[4]
Prove countable categoricity of dense linear orders by back-and-forth.
\end{exercise}

\begin{exercise}[4]
Study the Ryll--Nardzewski theorem conceptually.
\end{exercise}

\begin{exercise}[4]
Study quantifier elimination for dense linear orders.
\end{exercise}

\begin{exercise}[4]
Study quantifier elimination for algebraically closed fields.
\end{exercise}

\begin{exercise}[4]
Study quantifier elimination for real closed fields.
\end{exercise}

\begin{exercise}[4]
Study the Tarski--Seidenberg theorem.
\end{exercise}

\begin{exercise}[4]
Study Łoś's theorem in detail.
\end{exercise}

\begin{exercise}[4]
Construct an ultrapower of a simple structure.
\end{exercise}

\begin{exercise}[4]
Study nonprincipal ultrafilters and nonstandard arithmetic.
\end{exercise}

\begin{exercise}[4]
Study the model theory of modules.
\end{exercise}

\begin{exercise}[4]
Study o-minimal structures.
\end{exercise}

\begin{exercise}[4]
Study stability and the order property.
\end{exercise}

\begin{exercise}[4]
Study Morley's categoricity theorem conceptually.
\end{exercise}

\begin{exercise}[4]
Study Fraïssé classes and Fraïssé limits.
\end{exercise}

%%%%%%%%%%%%%%%%%%%%%%%%%%%%%%%%%%%%%%%%%%%%%%%%%%%%%%%%%%%%
\subsection{Conceptual Summary}
\label{subsec:model-theory-summary}
%%%%%%%%%%%%%%%%%%%%%%%%%%%%%%%%%%%%%%%%%%%%%%%%%%%%%%%%%%%%

Model theory studies the relationship between

\[
\boxed{
\text{formal theories}
}
\]

and

\[
\boxed{
\text{mathematical structures}.
}
\]

For a structure \(\mathcal{M}\),

\[
\boxed{
\operatorname{Th}(\mathcal{M})
}
\]

records all first-order sentences true in that structure.

Two structures are elementarily equivalent when

\[
\boxed{
\mathcal{M}\equiv\mathcal{N}.
}
\]

An elementary substructure satisfies

\[
\boxed{
\mathcal{M}\preccurlyeq\mathcal{N},
}
\]

meaning that every first-order formula with parameters from \(M\) has the
same truth value in both structures.

Definability studies subsets and functions described by formulas.

Compactness states that

\[
\boxed{
T
\text{ satisfiable}
\iff
\text{every finite subset of }T
\text{ satisfiable}.
}
\]

The Löwenheim--Skolem theorems show that infinite first-order theories often
have models of many cardinalities.

Types describe possible first-order properties of elements, while
saturated models realize large collections of consistent types.

Ultraproducts provide new structures satisfying first-order properties
according to Łoś's theorem.

Model theory therefore connects:

\[
\boxed{
\text{logic}
+
\text{algebra}
+
\text{geometry}
+
\text{cardinality}
+
\text{definability}
+
\text{classification}.
}
\]

%%%%%%%%%%%%%%%%%%%%%%%%%%%%%%%%%%%%%%%%%%%%%%%%%%%%%%%%%%%%
\subsection{Section Connections}
\label{subsec:model-theory-connections}
%%%%%%%%%%%%%%%%%%%%%%%%%%%%%%%%%%%%%%%%%%%%%%%%%%%%%%%%%%%%

\chapterconnection{
Model Theory develops the semantic side of the formal framework introduced
in Section~\ref{sec:formal-logic}. Axiomatic Set Theory provides an
important source of first-order theories and models, while Model Theory
studies structures through elementary equivalence, definability,
compactness, Löwenheim--Skolem, types, categoricity, ultraproducts, and
classification. These ideas explain both the expressive power and the
limitations of first-order axiomatization. The next section, Proof Theory,
shifts from semantic structures to the syntactic analysis of formal
derivations, proof systems, consistency, normalization, cut elimination,
and proof-theoretic strength.
}

%%%%%%%%%%%%%%%%%%%%%%%%%%%%%%%%%%%%%%%%%%%%%%%%%%%%%%%%%%%%
% SECTION SOURCE TRACKING
%%%%%%%%%%%%%%%%%%%%%%%%%%%%%%%%%%%%%%%%%%%%%%%%%%%%%%%%%%%%

%------------------------------------------------------------
% 12.3 Model Theory
%------------------------------------------------------------
%
% Source basis:
%   - Standard first-order model theory.
%   - Standard compactness, Löwenheim--Skolem, definability,
%     ultraproduct, categoricity, and classification concepts.
%
% Used for:
%   - languages, theories, and structures;
%   - classes of models;
%   - theory of a structure;
%   - elementary equivalence;
%   - isomorphisms;
%   - embeddings and substructures;
%   - elementary embeddings;
%   - elementary substructures;
%   - Tarski--Vaught criterion;
%   - definable sets and functions;
%   - definability with parameters;
%   - definable closure;
%   - model-theoretic algebraic closure;
%   - complete theories;
%   - theory completeness versus logical completeness;
%   - partial and complete types;
%   - realization and omission of types;
%   - compactness;
%   - finite satisfiability;
%   - infinite model construction;
%   - nonstandard arithmetic;
%   - elementary extensions;
%   - elementary diagrams;
%   - downward and upward Löwenheim--Skolem;
%   - Skolem paradox;
%   - internal and external cardinality;
%   - categoricity;
%   - dense linear orders;
%   - back-and-forth;
%   - quantifier elimination;
%   - model completeness;
%   - prime models;
%   - saturated models;
%   - homogeneity;
%   - automorphisms;
%   - ultrafilters;
%   - ultraproducts;
%   - ultrapowers;
%   - Łoś's theorem;
%   - nonstandard analysis preview;
%   - infinitesimals;
%   - algebraically closed fields;
%   - real closed fields;
%   - semialgebraic sets;
%   - Tarski--Seidenberg principle;
%   - o-minimality;
%   - finite model theory;
%   - stability theory preview;
%   - order property;
%   - classification theory;
%   - model-theoretic independence;
%   - types and automorphisms;
%   - Fraïssé theory preview;
%   - amalgamation;
%   - random graph preview;
%   - applications and exercises.
%
% Notes:
%   - Proof Theory is developed next in Section 12.4.
%   - Formal Logic appears in Section 12.1.
%   - Axiomatic Set Theory appears in Section 12.2.
%   - Computability Theory follows in Section 12.5.
%
%%%%%%%%%%%%%%%%%%%%%%%%%%%%%%%%%%%%%%%%%%%%%%%%%%%%%%%%%%%%